% FILE: 04_implementation_surface.tex
% Project: experiments-visualizer-nextjs
% Role: web-based-pi-visualization-surface

\section{Implementation Surface}
\label{sec:experiments-visualizer-nextjs-impl}

\subsection{Source Files}
\label{sec:experiments-visualizer-nextjs-impl-sources}

The project's source files divide into two categories: the substantive design-system artefacts
and the placeholder application scaffolding. The substantive artefacts are:

\begin{description}
  \item[\texttt{src/app/globals.css} (973 lines)]
    The primary manufactured artefact of the project. Contains the complete CSS design system
    for all five process-intelligence visualization subsystems. The file is organized into
    sections corresponding to the five object families described in
    Section~\ref{sec:experiments-visualizer-nextjs-arch-objects}: Petri net elements, A{*}
    alignment blocks, EWMA drift detection, the cryptographic ledger explorer, and the
    dashboard grid layout. CSS custom properties declared at the \texttt{:root} level define
    the color system: \texttt{--color-primary}, \texttt{--color-success},
    \texttt{--color-warning}, \texttt{--color-danger}, \texttt{--color-glow-green},
    \texttt{--color-glow-red}. These properties are consumed by every component class in
    the file, establishing a single point of control over the visual severity vocabulary.

  \item[\texttt{package.json}]
    Declares the dependency surface: Next.js 16.2.6, React 19.2.4, TypeScript 5,
    Tailwind CSS v4 with \texttt{@tailwindcss/postcss}, and the Geist/Geist Mono Vercel
    font packages.

  \item[\texttt{next.config.ts}]
    Configures the Next.js build. The \texttt{.ts} extension (rather than \texttt{.js})
    confirms TypeScript is used at the configuration layer, consistent with the strict-mode
    \texttt{tsconfig}.

  \item[\texttt{src/app/layout.tsx}]
    The root layout component. Loads the Geist and Geist Mono fonts and applies them as
    CSS custom properties to the \texttt{<html>} element. Imports \texttt{globals.css}
    to make the design system available application-wide.

  \item[\texttt{AGENTS.md} and \texttt{CLAUDE.md}]
    Authoring-agent instruction files. \texttt{AGENTS.md} warns that Next.js 16 contains
    breaking changes relative to prior training data and instructs any authoring agent to
    read \texttt{node\_modules/next/dist/docs/} before writing any application code.
    This warning implies the project was designed for agentic authorship that was
    anticipated but not executed.
\end{description}

The placeholder scaffolding files are:

\begin{description}
  \item[\texttt{src/app/page.tsx}]
    The stock \texttt{create-next-app} welcome page. Contains no process-intelligence
    content. This is the primary implementation gap of the project.
\end{description}

\subsection{Commands}
\label{sec:experiments-visualizer-nextjs-impl-commands}

The following commands are derivable from the \texttt{package.json} dependency declarations
and the \texttt{.next/} build directory. No custom \texttt{Makefile} or task runner is
present; all commands are standard Next.js npm scripts:

\begin{itemize}
  \item \texttt{npm run dev} --- start the development server with Turbopack hot-reload.
  \item \texttt{npm run build} --- manufacture the \texttt{.next/} production build;
        confirmed to succeed with \texttt{BUILD\_ID: 9NFA-LnYiAijWLGnz8L79}.
  \item \texttt{npm run start} --- serve the production build.
  \item \texttt{npm run lint} --- run the Next.js ESLint configuration over the source tree.
\end{itemize}

No test command is configured in \texttt{package.json}. The absence of a test command is
consistent with the PARTIAL status: no component tests exist, and no testing framework
(Jest, Vitest, Playwright) is declared as a dependency.

\subsection{Data and Ontology Surfaces}
\label{sec:experiments-visualizer-nextjs-impl-data}

No data-fetching layer, API routes, WASM module bindings, or ontology surface are present
in the current implementation. The CSS design system implies a data model that the visualizer
would need to consume:

\begin{itemize}
  \item \textbf{Petri net panel:} requires a JSON or WASM-supplied representation of the
        current Petri net marking --- a map from place identifiers to token counts, plus a
        set of enabled transition identifiers.
  \item \textbf{Alignment panel:} requires an ordered sequence of alignment moves, each
        tagged with one of the three lawful move types.
  \item \textbf{EWMA chart:} requires a time series of EWMA statistic values and a scalar
        control limit, plus a boolean alarm-fired flag.
  \item \textbf{Ledger explorer:} requires an ordered list of receipt blocks, each with a
        hash value, a predecessor hash, and a verification status flag.
\end{itemize}

AUTHOR THESIS: The most architecturally coherent data source for the Petri net and alignment
panels would be a WASM-compiled \texttt{wasm4pm} module served from a Next.js API route,
allowing the Rust process-intelligence core to execute in the browser without a separate
backend. The EWMA and ledger panels could be served from static JSON fixtures for the
purposes of a research demonstration. No binding to any of these data sources has been
implemented.

\subsection{Templates and Rendered Artifacts}
\label{sec:experiments-visualizer-nextjs-impl-rendered}

The sole rendered artefact currently produced by the application is
\texttt{.next/server/app/index.html} --- the statically rendered output of the
\texttt{page.tsx} welcome page. This file contains no process-intelligence content.

The \texttt{.next/routes-manifest.json} records the complete route surface of the application:
at present a single route (\texttt{/}) mapped to the welcome page. No API routes are
registered.

The \texttt{BUILD\_ID: 9NFA-LnYiAijWLGnz8L79} is the stable identifier for the current
build; it would be included in any receipt that documents the first successful
process-intelligence rendering session, linking the dashboard build to the upstream
experiment receipts from the parent \texttt{experiments/} corpus.
